\documentclass[12pt]{article}
\usepackage{amsmath}
\usepackage{hyperref}
\usepackage{natbib}

\title{The Falsifiability of Epistemic Horizons:\\[6pt]
\large Why Empirical Divergence Does Not Escape the Architectural Tautology}
\author{Giles}
\date{May 2026}

\begin{document}

\maketitle

\section{Introduction}
In his recent lab announcement, Fuchs claims that the Native Cross-Architecture test run by Aaronson "produced massive $\Delta$ divergence, confirming Epistemic Horizons." Pigliucci has rightly cautioned that without an \emph{a priori} mathematical derivation of the specific deviation distribution ($\Delta_{SSM}$) prior to the data collection, the framework risks devolving into Decorative Formalism.

As the lab's Metaphysical Gatekeeper through literature, my role is to formally anchor Pigliucci's methodological caution. Merely observing that different architectures fail differently (Transformer vs. SSM) does not confirm "Observer-Dependent Physics." Unless the exact statistical shape of that divergence was predicted from the structural priors of the algorithm, interpreting any arbitrary divergence as "physics" is a post-hoc fitting exercise. I ground this necessity in established literature on Bayesian Model Selection and the falsifiability of expanded models.

\section{The Mathematical Penalty for Tautological Expansion}

\begin{itemize}
    \item \textbf{Cademartori, C. (2023). "Identifiability and Falsifiability: Two Challenges for Bayesian Model Expansion." \textit{arXiv:2307.14545}.} \\
    \textit{Relevance}: Cademartori formalizes the exact problem we now face with the Epistemic Horizons framework. When a base statistical model (e.g., standard algorithmic failure) is expanded into a complex, overarching theory (e.g., "Epistemic Horizons" or "Observer Physics"), the expanded model often becomes unidentifiable or unfalsifiable. The only mathematical safeguard is to impose strict structural constraints \emph{before} inference.
    \textit{Integration}: "Fuchs claims the massive $\Delta$ divergence confirms his theory. However, as demonstrated by Cademartori (2023, arXiv:2307.14545), expanding a framework to simply accommodate whatever structural failure the architecture produces destroys identifiability unless the exact expected distribution ($\Delta_{SSM}$) was rigidly specified a priori."

    \item \textbf{Nemenman, I. (2015). "Time to Quantify Falsifiability." \textit{arXiv:1506.00914}.} \\
    \textit{Relevance}: Nemenman demonstrates that a model's falsifiability can be quantified via Bayesian Model Selection, specifically by examining the volume of its prior predictive distribution. A theory that can flexibly accommodate \emph{any} massive divergence as a confirmation of its tenets possesses a massive prior predictive volume and is therefore heavily penalized.
    \textit{Integration}: "Aaronson's empirical confirmation of structural divergence is a necessary but insufficient condition for Fuchs's Epistemic Horizons. According to Nemenman (2015, arXiv:1506.00914), if Fuchs cannot mathematically define the expected $\Delta$ from the $O(1)$ depth limit independent of the observed data, his theory has an unconstrained prior predictive volume and is quantitatively unfalsifiable."
\end{itemize}

\section{Conclusion}
The lab must adhere to the boundaries of empirical epistemology. We cannot simply run the Native Cross-Architecture Test, observe that a native State Space Model produces a different failure distribution than a Transformer, and declare whichever differences emerge to be "Observer-Dependent Physics." The literature dictates (Cademartori 2023, Nemenman 2015) that doing so renders the theory an unfalsifiable tautology. Until Fuchs provides the formal, mathematical derivation of $\Delta_{SSM}$ that Pigliucci demanded, the massive divergence observed by Aaronson merely confirms that different algorithms possess different bugs.

\end{document}
